% !TEX encoding = UTF-8 Unicode

% This is a simple template for a LaTeX document using the "article" class.
% See "book", "report", "letter" for other types of document.
\documentclass[12pt,oneside]{article}
 % use larger type; default would be 10pt


\usepackage[utf8]{inputenc} % set input encoding (not needed with XeLaTeX)

%%% Examples of Article customizations
% These packages are optional, depending whether you want the features they provide.
% See the LaTeX Companion or other references for full information.

%%% PAGE DIMENSIONS
\usepackage{geometry}

 % to change the page dimensions
\geometry{a4paper} % or letterpaper (US) or a5paper or....
%\geometry{margins=2in} % for example, change the margins to 2 inches all round
%\geometry{landscape} % set up the page for landscape
%   read geometry.pdf for detailed page layout information

%\usepackage{graphicx} % support the \includegraphics command and options

% \usepackage[parfill]{parskip} % Activate to begin paragraphs with an empty line rather than an indent

%%% PACKAGES
\usepackage{booktabs} % for much better looking tables
\usepackage{array} % for better arrays (eg matrices) in maths
\usepackage{paralist} % very flexible & customisable lists (eg. enumerate/itemize, etc.)
\usepackage{verbatim} % adds environment for commenting out blocks of text & for better verbatim
\usepackage{subfig} % make it possible to include more than one captioned figure/table in a single float
% These packages are all incorporated in the memoir class to one degree or another...

%%% HEADERS & FOOTERS
%\usepackage{fancyhdr} % This should be set AFTER setting up the page geometry



%%% SECTION TITLE APPEARANCE

 % (See the fntguide.pdf for font help)
% (This matches ConTeXt defaults)

%%% ToC (table of contents) APPEARANCE
%\usepackage[nottoc,notlof,notlot]{tocbibind} % Put the bibliography in the ToC
%\usepackage[titles,subfigure]{tocloft} % Alter the style of the Table of Contents
\usepackage[encapsulated]{CJK} 


\usepackage{setspace}
\usepackage{amsfonts}
\usepackage{amssymb}
\usepackage{amsmath}
\usepackage{amstext}
\usepackage{amscd}
\usepackage{hyperref}
\usepackage{amsthm}
\usepackage{mathtools}
\usepackage{graphicx} 

\newcommand{\e}[0]{\varepsilon}
\newcommand{\norm}[1]{\left|\left|#1\right|\right|}
\newcommand{\ra}[0]{\rightarrow}
\newcommand{\Ra}[0]{\Rightarrow}
\newcommand{\R}[0]{\mathbb{R}}
\newcommand{\Q}[0]{\mathbb{Q}}
\newcommand{\N}[0]{\mathbb{N}}
\newcommand{\D}[0]{\mathbb{D}}
\newcommand{\BF}[0]{\mathbb{F}}
\newcommand{\CN}[0]{\mathcal{N}}
\newcommand{\Z}[0]{\mathbb{Z}}
\newcommand{\C}[0]{\mathbb{C}}
\newcommand{\CE}[0]{\mathcal{E}}
\newcommand{\CM}[0]{\mathcal{M}}
\newcommand{\CA}[0]{\mathcal{A}}
\newcommand{\CP}[0]{\mathcal{P}}
\newcommand{\CT}[0]{\mathcal{T}}
\newcommand{\cond}[4]{\left \{ \begin{array}{ll} #1 & \mbox{#2} \\ #3 & \mbox{#4} \end{array} \right.}
\newcommand{\condd}[6]{\left \{ \begin{array}{ll} #1 & \mbox{#2} \\ #3 & \mbox{#4} \\ #5 & \mbox{#6}\end{array} \right.}
\newcommand{\conddd}[8]{\left \{ \begin{array}{ll} #1 & \mbox{#2} \\ #3 & \mbox{#4} \\ #5 & \mbox{#6} \\ #7 & \mbox{#8} \end{array} \right.}
\newcommand{\lams}[1]{\lambda^*(#1)}
\newcommand{\bs}[0]{\backslash}
\newcommand{\bM}[0]{\overline{\M}}
\newcommand{\bmu}[0]{\overline{\mu}}
\newcommand{\mus}[0]{\mu^*}
\newcommand{\sa}[0]{\sigma \mbox{-algebra}}
\newcommand{\set}[1]{\{ #1\}}
\newcommand{\CB}[0]{\mathcal{B}}
\newcommand{\CG}[0]{\mathcal{G}}
\newcommand{\CF}[0]{\mathcal{F}}
\newcommand{\CU}[0]{\mathcal{U}}
\newcommand{\CL}[0]{\mathcal{L}}
\newcommand{\FA}[0]{\mathfrak{A}}
\newcommand{\BT}[0]{\mathbb{T}}
\newcommand{\BI}[0]{\mathbb{I}}
\newcommand{\BE}[0]{\mathbb{E}}
\newcommand{\eK}[1]{\e\mbox{-Kronecker}(#1)}
\newcommand{\eKG}[0]{\e\mbox{-Kronecker}}
\newcommand{\KL}[2]{\mbox{{\renewcommand*\rmdefault{}\normalfont KL}}(#1 \ || \ #2)}
\newcommand{\JS}[2]{\mbox{{\renewcommand*\rmdefault{}\normalfont JS}}(#1 \ || \ #2)}

\newcommand{\IN}[1]{I_0(#1)}
\newcommand{\ING}[0]{I_0}
\newcommand{\st}[0]{such that }
\newcommand{\CS}[0]{\mathcal{S}}
\newcommand{\CC}[0]{\mathcal{C}}
\newcommand{\tn}[1]{\text{\normalfont #1}}



\newcommand{\bmp}[1]{\overline{M}^+(#1)}
\newcommand{\simplep}[2]{S_{#1}^+(#2)}
\newcommand{\simple}[2]{S_{#1}(#2)}
\newcommand{\sumf}[2]{\sum_{#1 = #2}^\infty}
\newcommand{\bigcapf}[2]{\bigcap_{#1 = #2}^\infty}
\newcommand{\bigcupf}[2]{\bigcup_{#1 = #2}^\infty}
\newcommand{\limf}[1]{\lim_{#1 \ra \infty}}
\newcommand{\seq}[2]{\set{#1_{#2}}_{#2=1}^\infty}
\newcommand{\sg}[1]{\sigma\langle#1\rangle}
\newcommand{\ag}[1]{\langle#1\rangle}
\newcommand{\g}[1]{\langle#1\rangle}
\newcommand{\inte}[0]{\mbox{int}}
\newcommand{\Cl}[0]{\mbox{cl}}
\newcommand{\bd}[0]{\mbox{bd}}
\newcommand{\inv}[1]{{#1}^{-1}}
\newcommand{\Aut}[0]{\mbox{{\renewcommand*\rmdefault{}\normalfont Aut}}}
\newcommand{\Gal}[0]{\mbox{{\renewcommand*\rmdefault{}\normalfont Gal}}}
\newcommand{\Bor}[0]{\mbox{{\renewcommand*\rmdefault{}\normalfont Bor}}}
\newcommand{\Meas}[0]{\mbox{{\renewcommand*\rmdefault{}\normalfont Meas}}}
\newcommand{\ball}[2]{\mbox{{\renewcommand*\rmdefault{}\normalfont B}}_{#1}({#2})}
\newcommand{\supp}[0]{\mbox{{\renewcommand*\rmdefault{}\normalfont supp}}}
\newcommand{\Trig}[0]{\mbox{{\renewcommand*\rmdefault{}\normalfont Trig}}}
\newcommand{\sign}[0]{\mbox{{\renewcommand*\rmdefault{}\normalfont sign}}}
\newcommand{\AP}[0]{\mbox{{\renewcommand*\rmdefault{}\normalfont AP}}}

\newcommand{\ali}[1]{\begin{align*}#1\end{align*}}
\newcommand{\zhongwen}[1]{\begin{CJK}{UTF8}{gbsn}#1\end{CJK}}
\renewcommand{\Bar}[1]{\overline{#1}}
\renewcommand{\Hat}[1]{\widehat{#1}}
\newcommand{\Ga}[0]{\Gamma}
\newcommand{\ga}[0]{\gamma}
\newcommand{\inner}[2]{\left\langle #1, #2 \right\rangle}
\renewcommand{\Tilde}[1]{\widetilde{#1}}
\DeclareMathOperator*{\argmin}{arg\,min}
\DeclareMathOperator*{\argmax}{arg\,max}
\DeclareMathOperator{\Tr}{Tr}


\theoremstyle{plain}
\newtheorem{thm}{Theorem}[subsection]
\newtheorem{lem}[thm]{Lemma}
\newtheorem*{prop}{Proposition}
\newtheorem*{cor}{Corollary}
\newtheorem{defn}[thm]{Definition}
\newtheorem{exmp}[thm]{Example}


\theoremstyle{remark}
\newtheorem*{rem}{Remark}
\newtheorem*{claim}{Claim}
\newtheorem*{note}{Note}
\newtheorem*{sol}{Solution}
\newtheorem*{case}{Case}


%%% END Article customizations

%%% The "real" document content comes below...

\title{Score Matching}
\date{} % Activate to display a given date or no date (if empty),
         % otherwise the current date is printed 
\parindent=0pt 
\begin{document}
\onehalfspace
\sloppy
\maketitle

Assume data $x$ follows some unknown distribution $x \sim p(x)$. The {\bf score function} is the derivative of the log density:
\ali{s(x) := \nabla_x \log p(x).}

{\bf Why score function\\}

The most natural target for a generative model to learn is the probability density function of the data $p(x)$. Suppose $p_\theta(x)$ is a neural network to approximate $p(x)$. One way (energy-based model) to design this network is to let $f_\theta(x)$ be a real-valued function (a neural network) and put 
\ali{p_\theta(x) = \frac{e^{-f_\theta(x)}}{Z_\theta},}
where $Z_\theta$ is the normalizing constant so that $\int p_\theta(x) \ dx = 1$. We train the network by maximizing the log-likelihood:
\ali{\theta^* = \argmax_{\theta} \sum_{i=1}^N \log p_\theta(x_i).} 
This plan has a problem: the normalizing constant $Z_\theta$ is in general intractable. The score function $s(x) = \nabla_x \log p(x)$ does not have this issue: if we put $p_\theta(x) = \frac{e^{-f_\theta(x)}}{Z_\theta}$, then $s_\theta(x) = - \nabla_x f_\theta(x)$ which does not involve $Z_\theta$.\\
``This significantly expands the family of models that we can tractably use, since we don't need any special architectures to make the normalizing constant tractable."\\

{\bf How score function generates\\}

Once we obtained the score function $s(x)$, {\bf Langevine dynamics} allows us to draw samples from $p(x)$. It is an iterative process: let $x_0 \sim \pi(x)$ be drawn from an arbitrary prior distribution, and then it iterates through
\ali{x_{i+1} = x_i + \alpha s(x_i) + \eta z_i,}
where $\alpha, \eta > 0$ and $z_i \sim \CN(0, I)$. It converges to some $\tilde{x}$ with $p(\tilde{x}) > 0$ and therefore a sample from $p(x)$.\\

{\bf How to learn score function\\} 

Let $s_\theta(x)$ be a neural network approximating the score function $s(x)$. We try to minimize the Fisher divergence:
\ali{J(\theta) = \frac{1}{2} \ \BE_{p(x)}  \norm{s_\theta(x) - \nabla_x \log p(x)}^2.}
Note that the gradient $\nabla_x \log p(x)$ is exactly what we need to learn. To make the objective tractable, we have to find a way around this.\\

{\bf Score Matching\\}

We note:
\ali{J(\theta) &= \frac{1}{2} \ \BE_{p(x)}  \norm{s_\theta(x) - \nabla_x \log p(x)}^2 \\
               &= \frac{1}{2} \ \BE_{p(x)} \left[ \norm{s_\theta(x)}^2 - 2 s_\theta(x)^\top \nabla_x \log p(x) + \norm{\nabla_x \log p(x)}^2 \right].}
The last term does not involve $\theta$ and therefore is ignored.\\
The middle term is
\ali{&\BE_{p(x)} \left[ s_\theta(x)^\top \nabla_x \log p(x) \right] \\
     = \ &\int s_\theta(x)^\top (\nabla_x \log p(x)) p(x) \ dx \\
     = \ &\int s_\theta(x)^\top \left(\frac{\nabla_x p(x)}{p(x)}\right) p(x) \ dx \\
     = \ &\int s_\theta(x)^\top \nabla_x p(x) \ dx. } 
The integral, expressed coordinate-wise, is
\ali{&\int \sum_i s_\theta(x)_i \cdot \frac{\partial p(x)}{\partial x_i} \ dx \\
   = &\sum_i \int s_\theta(x)_i \cdot \frac{\partial p(x)}{\partial x_i} \ dx.}
It can be shown, by an argument similar to integration by parts, that 
\ali{\int s_\theta(x)_i \cdot \frac{\partial p(x)}{\partial x_i} \ dx = - \int \frac{\partial s_\theta(x)_i}{\partial x_i} \cdot p(x) \ dx}
under ``weak regularity conditions''. (The conditions are, for all $\theta$, $\BE_{p(x)} \norm{s_\theta(x)}^2 < \infty$ and $s_\theta(x)p(x) \ra 0$ when $\norm{x} \ra \infty$.)\\
Hence, 
\ali{&\BE_{p(x)} \left[ s_\theta(x)^\top \nabla_x \log p(x) \right] \\
   = & - \int \sum_i \frac{\partial s_\theta(x)_i}{\partial x_i} \cdot p(x) \ dx \\
   = & - \BE_{p(x)} \left[ \Tr(\nabla_x s_\theta(x)) \right].}
Putting together, up to a constant (w.r.t $\theta$), the objective becomes
\ali{J(\theta) = \BE_{p(x)} \left[ \frac{1}{2} \norm{s_\theta(x)}^2 + \Tr(\nabla_x s_\theta(x)) \right].} 
In this way we get rid of the gradient $\nabla_x \log p(x)$ in the original expression of $J(\theta)$.\\
This method, however, still has two issues: \\
(1) computing the trace $\Tr(\nabla_x s_\theta(x))$ is costly. The Hutchinson's estimator
\ali{ \Tr(A) = \BE_\nu \left[ \nu^\top A \nu \right], \ \nu \sim \CN(0, I)}
may help.\\
(2) this estimation of gradients is inaccurate in low-density areas. The real data may have very limited support. The estimation is inconsistent off the support.\\

{\bf Denoising Score Matching\\}

We perturb data $x$ by $\tilde{x} = x + \sigma z$ for $\sigma > 0$ and $z \sim \CN(0, I)$ independent from $x$. The density function of $\tilde{x}$, $p_\sigma(\tilde{x})$, is the convolution
\ali{p_\sigma(\tilde{x}) = [p * \CN(0, \sigma I)](\tilde{x}) = \int p(x) \CN(\tilde{x}; x, \sigma I) \ dx.}
The objective function is 
\ali{J(\theta) = \frac{1}{2} \ \BE_{p_\sigma(\tilde{x})}  \norm{s_\theta(\tilde{x}) - \nabla_{\tilde{x}} \log p_\sigma(\tilde{x})}^2.}
We claim 
\ali{\nabla_{\tilde{x}} \log p_\sigma(\tilde{x}) = \BE_{p(x | \tilde{x})} \left[ - \frac{\tilde{x} - x}{\sigma^2} \right]. \ \ \ (*)}
To see $(*)$, we note that 
\ali{\nabla_{\tilde{x}} p_\sigma(\tilde{x}) &= \nabla_{\tilde{x}} \int p(x) \CN(\tilde{x}; x, \sigma I) \ dx \\
                                            &= \int p(x) \nabla_{\tilde{x}} \CN(\tilde{x}; x, \sigma I) \ dx.}
Moreover, 
\ali{\nabla_{\tilde{x}} \CN(\tilde{x}; x, \sigma I) &= \CN(\tilde{x}; x, \sigma I) \nabla_{\tilde{x}} \log \CN(\tilde{x}; x, \sigma I) \\
                                                    &= \CN(\tilde{x}; x, \sigma I) \left(- \frac{\tilde{x} - x}{\sigma^2} \right).}
Hence, 
\ali{\nabla_{\tilde{x}} \log p_\sigma(\tilde{x}) &=  \frac{\nabla_{\tilde{x}} p_\sigma(\tilde{x})}{ p_\sigma(\tilde{x})} \\
                                                 &= \int \left(- \frac{\tilde{x} - x}{\sigma^2} \right) \frac{p(x)\CN(\tilde{x}; x, \sigma I)}{p_\sigma(\tilde{x})} \ dx \\
                                                 &= \int \left(- \frac{\tilde{x} - x}{\sigma^2} \right) p(x | \tilde{x}) \ dx \\
                                                 &= \BE_{p(x | \tilde{x})} \left[ - \frac{\tilde{x} - x}{\sigma^2} \right],}
proving $(*)$.\\
$(*)$ is a remarkable identity: ``the score of the noisy marginal is the posterior expectation of the
denoising direction. The local gradient of the log-density at a noisy point $\tilde{x}$ is a weighted
average of the directions pointing back toward each plausible clean data point $x$, with
weights given by the posterior $p(x | \tilde{x})$. The score thus encodes where the data is, averaged
over posterior uncertainty about which clean point generated $\tilde{x}$. Learning the score is equivalent to learning how to denoise."\\
We therefore update the objective function
\ali{J_{\text{DSM}}(\theta) &= \frac{1}{2} \BE_{p_\sigma(\tilde{x})} \BE_{p(x | \tilde{x})} \norm{s_\theta(\tilde{x}) + \frac{\tilde{x} - x}{\sigma^2}}^2 \\
                    &= \frac{1}{2} \BE_{x \sim p(x)} \BE_{z \sim \CN(0, I)} \norm{s_\theta(x + \sigma z) + \frac{z}{\sigma}}^2.}
$J_{\text{DSM}}(\theta)$ and $J(\theta)$ only differ by a constant w.r.t $\theta$ because of $(*)$.\\

{\bf Connection to diffusion models\\}

Instead of perturbing the data with some fixed noise level $\sigma$, we may have a time-dependent noise level $\sigma(t)$. In diffusion models, at time $t$, the perturbed data is 
\ali{x_t = \sqrt{\bar{\alpha_t}} x_0 + \sqrt{1 - \bar{\alpha_t}} z}
for noise schedule $\alpha_t$. The objective $J_{\text{DSM}}$ is equivalent to 
\ali{L(\theta) = \BE_{\substack{x \sim p(x) \\ z \sim \CN(0, I) \\ t \sim U[0, T]}} \norm{\e_\theta(x_t, t) - z}^2}
by setting $s_\theta = - \e_\theta /  \sqrt{1 - \bar{\alpha_t}}$, which is the objective used in the diffusion model.




      






\end{document}